% ==============================================================================
% capitulos/04-tecnologias.tex - Selección de Tecnologías
% ==============================================================================

\chapter{Selección de Tecnologías}
\label{cap:tecnologias}

\section{Definición de arquitectura de software}
La arquitectura de software se entiende como el conjunto de patrones y abstracciones que proporcionan el marco de referencia necesario para llevar a cabo la construcción de un sistema. Funciona como el esqueleto del software, ya que en este se define cómo se organizan los componentes, cómo interactúan entre sí y bajo qué principios se rigen para cumplir con los objetivos del negocio [13]. 

\subsection{Importancia en el desarrollo de sistemas}
Su importancia radica en que permite a los equipos de desarrollo tener una visión clara y compartida de la solución, facilitando la toma de decisiones críticas antes de iniciar la etapa de codificación detallada. 

\subsection{Factores de calidad}
Para asegurar que la estructura seleccionada sea la adecuada, esta debe ser evaluada bajo atributos de calidad que garanticen su viabilidad a largo plazo. De acuerdo con los modelos de calidad internacionales como el ISO/IEC 9126, estos atributos permiten medir la capacidad del software para satisfacer necesidades explícitas e implícitas bajo condiciones determinadas [14]. En sistemas de gestión institucional, destacan los siguientes pilares:
\begin{itemize}
   \item \textbf{Rendimiento:}  Evalúa el comportamiento temporal del sistema y la utilización óptima de recursos al procesar transacciones y responder a las peticiones del usuario [14].
   \item \textbf{Mantenibilidad:}  El esfuerzo necesario para realizar modificaciones, incluyendo correcciones, mejoras o adaptaciones a nuevos requisitos normativos, garantizando la estabilidad del software [14].
    \item \textbf{Escalabilidad:} Se define como la capacidad de un sistema o aplicación para ampliarse o reducirse con el fin de ajustarse a los cambios en las necesidades de procesamiento y en el volumen de demanda. Un sistema escalable puede asimilar un aumento de la carga de trabajo, ya sea añadiendo recursos de hardware o actualizando el software, sin degradar su rendimiento ni interrumpir su funcionamiento [15].
\end{itemize}

\subsection{Clasificación de arquitecturas}

Históricamente, las arquitecturas se han clasificado según la forma en que distribuyen sus responsabilidades. De acuerdo con la literatura y los referentes de la industria, destacan las siguientes:

\begin{itemize}
    \item \textbf{Arquitecturas Monolíticas:} Modelo tradicional donde la interfaz, la lógica de negocio y el acceso a datos operan como una unidad indivisible y centralizada dentro de un mismo programa [16].
    \item \textbf{Arquitecturas Distribuidas:} Sistemas donde los componentes están alojados en distintos equipos de una red y colaboran mediante el intercambio de mensajes para compartir la carga de procesamiento [17].
    \item \textbf{Arquitecturas de Microservicios:} Enfoque que divide una aplicación en servicios pequeños, independientes y comunicados a través de interfaces (APIs), lo que permite escalar y actualizar funciones de manera aislada [18].
    \item \textbf{Arquitecturas por Capas (Tres Capas):} Patrón que organiza el software en niveles independientes (presentación, lógica y datos) para separar responsabilidades y facilitar su mantenimiento a nivel de código [19].
\end{itemize}

\section{Arquitectura del Sistema}
\subsection{Arquitectura de Microservicios (Macro-arquitectura)}
\subsection{Arquitectura Interna de los Servicios: Patrón de 3 Capas (Micro-arquitectura)}
\subsection{Servidor Web y Proxy Inverso (Nginx)}

\section{Tecnologías de Desarrollo Backend}
\subsection{Lenguaje de Programación: Java}
\subsection{Framework Empresarial: Spring Boot}

\section{Tecnologías de Procesamiento de Inteligencia Artificial}
\subsection{Lenguaje de Programación Científica: Python}

\section{Tecnologías de Desarrollo Frontend}
\subsection{Framework de Interfaz de Usuario: Angular}

\section{Gestión de Bases de Datos}
\subsection{Sistema de Gestión de Bases de Datos Relacional (RDBMS): PostgreSQL}

\section{Infraestructura y Despliegue}
\subsection{Plataforma de Contenerización: Docker}

\section{Control de Versiones y Trabajo Colaborativo}
\subsection{Repositorio y Flujos de Trabajo: GitHub}
